\documentclass[aspectratio=169]{beamer}

\usepackage[utf8]{inputenc}
\usepackage[T1]{fontenc}
\usepackage[brazil]{babel}
\usepackage{amsmath, amssymb}
\usepackage[expansion=false,protrusion=false]{microtype}
\usepackage{csquotes}
\usepackage{tikz}
\usepackage{booktabs}
\usepackage{array}

% Colors — paleta padrão (azul escuro + laranja)
\definecolor{mblack}{RGB}{20, 20, 20}
\definecolor{mprimary}{RGB}{8,139,188}     % azul escuro
\definecolor{maccent}{RGB}{97,159,32}     % laranja
\definecolor{mwhite}{RGB}{255, 255, 255}

% Theme
\usetheme{default}
\usecolortheme{default}
\setbeamercolor{normal text}{fg=mblack, bg=white}
\setbeamercolor{frametitle}{fg=white, bg=mprimary}
\setbeamercolor{title}{fg=white}
\setbeamercolor{author}{fg=white}
\setbeamercolor{date}{fg=white}
\setbeamercolor{itemize item}{fg=maccent}
\setbeamercolor{itemize subitem}{fg=mprimary}
\setbeamercolor{enumerate item}{fg=maccent}
\setbeamercolor{enumerate subitem}{fg=mprimary}
\setbeamercolor{block title}{fg=white, bg=mprimary}
\setbeamercolor{block body}{fg=mblack, bg=mprimary!8}
\setbeamercolor{block title alerted}{fg=white, bg=maccent}
\setbeamercolor{block body alerted}{fg=mblack, bg=maccent!10}
\setbeamercolor{alerted text}{fg=maccent}

\setbeamerfont{frametitle}{series=\bfseries, size=\large}
\setbeamerfont{title}{series=\bfseries, size=\LARGE}
\setbeamerfont{block title}{series=\bfseries}

\setbeamertemplate{navigation symbols}{}

% Footer
\setbeamertemplate{footline}{%
  \leavevmode%
  \hbox{%
    \begin{beamercolorbox}[wd=0.7\paperwidth, ht=2.25ex, dp=1ex, leftskip=0.3cm]{normal text}%
      \tiny\color{mblack!50} Introdução ao Curso -- DECOM/UFOP
    \end{beamercolorbox}%
    \begin{beamercolorbox}[wd=0.3\paperwidth, ht=2.25ex, dp=1ex, rightskip=0.3cm, right]{normal text}%
      \tiny\color{mblack!50}\insertframenumber{} / \inserttotalframenumber
    \end{beamercolorbox}%
  }%
}

% Frametitle with accent rule
\setbeamertemplate{frametitle}{%
  \nointerlineskip%
  \begin{beamercolorbox}[wd=\paperwidth, ht=0.55cm, dp=0.2cm, leftskip=0.5cm]{frametitle}%
    \usebeamerfont{frametitle}\insertframetitle%
  \end{beamercolorbox}%
  \vspace{-0.1cm}%
  \textcolor{maccent}{\rule{\paperwidth}{2pt}}%
}

% Title page: three-band layout
\setbeamercolor{titlebar top}{bg=mwhite}
\setbeamercolor{titlebar main}{bg=mprimary}
\setbeamercolor{titlebar bottom}{bg=maccent}

\setbeamertemplate{title page}{%
  \vspace{0pt}%
  \begin{beamercolorbox}[wd=\paperwidth, ht=0.55cm, dp=0pt]{titlebar top}\end{beamercolorbox}%
  \begin{beamercolorbox}[wd=\paperwidth, ht=3.8cm, dp=0pt, center]{titlebar main}%
    \vspace{0.6cm}
    {\usebeamerfont{title}\color{white}\inserttitle\par}%
    \vspace{0.35cm}{\color{white}\small\insertauthor\par}%
    \vspace{0.15cm}{\color{white}\footnotesize\insertdate\par}%
    \vspace{0.25cm}
  \end{beamercolorbox}%
  \begin{beamercolorbox}[wd=\paperwidth, ht=1.3cm, dp=0pt]{titlebar bottom}\end{beamercolorbox}%
}

\setlength{\itemsep}{4pt}

% Highlight commands
\newcommand{\impt}[1]{\textcolor{maccent}{\textbf{#1}}}
\newcommand{\ntc}[1]{\textcolor{mprimary}{#1}}

\title{Estrutura Acadêmica da UFOP}
\author{Departamento e Colegiado de Curso}
\date{\today}

\begin{document}

% ============================================================
\begin{frame}[plain]\titlepage\end{frame}

% ============================================================
\begin{frame}{Sumário}
  \tableofcontents
\end{frame}

% ============================================================
\section{Onde estamos na estrutura da UFOP}

\begin{frame}{A UFOP em dois grandes blocos}
  \begin{columns}[T,onlytextwidth]
    \begin{column}{0.48\textwidth}
      \begin{block}{Administração Central}
        \begin{itemize}
          \item Reitoria
          \item Conselho Universitário
          \item Conselho Curador
        \end{itemize}
      \end{block}
    \end{column}
    \begin{column}{0.48\textwidth}
      \begin{block}{Unidades Acadêmicas}
        \begin{itemize}
          \item ICEB, EM, IFAC, \ldots
          \item DECOM $\subset$ ICEB
        \end{itemize}
      \end{block}
    \end{column}
  \end{columns}

  \vspace{0.4cm}
   
  \ntc{É dentro das Unidades Acadêmicas que vivem os dois órgãos desta aula.}
\end{frame}

% ============================================================
\begin{frame}{Composição de uma Unidade Acadêmica}
  \begin{itemize}
    \item Conselho da Unidade
    \item Diretoria e vice-diretoria
    \item Cursos de graduação e pós-graduação
    \item \impt{Colegiados de curso}
    \item \impt{Departamentos} (ou equivalentes)
    \item Núcleos e órgãos complementares
  \end{itemize}

  \vspace{0.4cm}
   
  \begin{alertblock}{A distinção essencial}
    \textbf{Departamento:} onde o \emph{professor} trabalha.\\
    \textbf{Colegiado:} quem cuida da formação do \emph{aluno} de um curso.
  \end{alertblock}
\end{frame}

% ============================================================
\begin{frame}{Exemplo concreto}
  Um docente do \impt{DECOM} pode ministrar disciplinas para:
  \begin{itemize}
    \item BCC, BIA
    \item Engenharias
    \item Estatística
    \item Matemática
    \item Química
  \end{itemize}

  \vspace{0.3cm}
   
  \ntc{Lotação:} um único departamento (DECOM).\\
  \ntc{Representação:} potencialmente em vários colegiados.
\end{frame}

% ============================================================
\section{Departamento}

\begin{frame}{Departamento --- definição}
  \begin{block}{Estatuto da UFOP}
    Órgão de \impt{lotação de professores} (e de parte dos técnicos-administrativos), voltado para objetivos comuns de \impt{ensino, pesquisa e extensão} em uma área do conhecimento.
  \end{block}

  \vspace{0.4cm}

  Toda Unidade Acadêmica é composta por departamentos --- e todo departamento está vinculado a uma Unidade.
\end{frame}

% ============================================================
\begin{frame}{Por que um departamento existe?}
  Três justificativas:
  \begin{enumerate}
    \item Áreas do conhecimento que abrange
    \item Linhas de pesquisa e projetos pedagógicos sustentados
    \item Recursos materiais e humanos necessários
  \end{enumerate}

  \vspace{0.4cm}

  \begin{block}{O DECOM, em particular}
    Algoritmos, IA, sistemas, teoria da computação\ldots\\
    + pesquisa associada + pessoal lotado.
  \end{block}
\end{frame}

% ============================================================
\begin{frame}{Quem compõe a Assembleia do Departamento}
  \begin{itemize}
    \item \impt{Todos os docentes} lotados no departamento
    \item Representante(s) dos \impt{técnicos-administrativos}
    \item Representante(s) do \impt{corpo discente} dos cursos atendidos
  \end{itemize}

  \vspace{0.4cm}

  \ntc{Discentes e técnicos têm representação paritária entre si.}\\
  \ntc{Deliberações: maioria dos presentes.}
\end{frame}

% ============================================================
\begin{frame}{Chefia e funcionamento}
  \begin{columns}[T,onlytextwidth]
    \begin{column}{0.48\textwidth}
      \begin{block}{Chefia}
        \begin{itemize}
          \item Chefe e vice-chefe
          \item Mandato de \impt{2 anos}
          \item Permitida uma recondução
          \item Chefe preside a assembleia
        \end{itemize}
      \end{block}
    \end{column}
    \begin{column}{0.48\textwidth}
      \begin{block}{Reuniões}
        \begin{itemize}
          \item Mínimo \impt{2 por semestre}
          \item Câmara: possível com \impt{>20 docentes}
          \item Com câmara: $\geq$ 1 assembleia plena/semestre
        \end{itemize}
      \end{block}
    \end{column}
  \end{columns}
\end{frame}

% ============================================================
\begin{frame}{O que o departamento decide --- parte 1}
  \begin{enumerate}
    \item Planejar ensino, pesquisa e extensão
    \item Elaborar planos de trabalho e \impt{capacitação} dos docentes
    \item Atribuir \impt{encargos de ensino} \\
          \ntc{(quem dá qual disciplina em cada semestre)}
    \item Propor aos colegiados as \impt{ementas, programas e bibliografias} das suas disciplinas
  \end{enumerate}
\end{frame}

% ============================================================
\begin{frame}{O que o departamento decide --- parte 2}
  \begin{enumerate}
    \setcounter{enumi}{4}
    \item Propor \impt{contratação, substituição, afastamento e dispensa} de docentes
    \item Eleger \impt{representantes nos colegiados de curso}
    \item Elaborar o \impt{relatório anual} de atividades
    \item Propor (com 2/3) afastamento ou destituição do chefe
  \end{enumerate}

  \vspace{0.4cm}
  \begin{alertblock}{Resumo}
    Departamento responde por \impt{oferta} (quem ensina, o que ensina, com qual conteúdo) e \impt{corpo docente}.
  \end{alertblock}
\end{frame}

% ============================================================
\begin{frame}{Atenção: ementa é \emph{proposta}, não \emph{aprovada}}
  \begin{block}{Direção da proposta}
    O DECOM \impt{propõe} a ementa de BCC501.\\
    O colegiado do BCC (e demais que usam a disciplina) \impt{aprova}.
  \end{block}

  \vspace{0.4cm}
   
  \ntc{Departamento define o conteúdo técnico; colegiado avalia no contexto do projeto pedagógico do curso.}
\end{frame}

% ============================================================
\section{Colegiado de Curso}

\begin{frame}{Colegiado --- definição}
  \begin{block}{Função}
    Órgão responsável pela \impt{coordenação didática} dos componentes curriculares do \impt{projeto pedagógico} de um curso.
  \end{block}

  \vspace{0.4cm}
   
  Cada curso de graduação e cada programa de pós tem \impt{o seu próprio colegiado}.

  \vspace{0.3cm}
   
  \ntc{Mudança de foco: o departamento se organiza por área; o colegiado, por curso.}
\end{frame}

% ============================================================
\begin{frame}{Quem compõe o colegiado}
  \begin{itemize}
    \item \impt{Representantes docentes} dos departamentos que oferecem disciplinas do curso
    \item \impt{Representação estudantil} (mandato de 1 ano, uma recondução)
  \end{itemize}

  \vspace{0.4cm}
   
  \begin{block}{Regra de proporcionalidade}
    1 representante a cada \impt{180h} ofertadas pelo departamento.\\
    Máximo de \impt{5 membros} por departamento.
  \end{block}
\end{frame}

% ============================================================
\begin{frame}{Coordenação e NDE}
  \begin{columns}[T,onlytextwidth]
    \begin{column}{0.48\textwidth}
      \begin{block}{Coordenação}
        \begin{itemize}
          \item Coordenador e vice (docentes)
          \item Mandato de \impt{2 anos}
          \item Uma recondução
          \item Coordenador preside o colegiado
        \end{itemize}
      \end{block}
    \end{column}
    \begin{column}{0.48\textwidth}
      \begin{block}{NDE}
        \begin{itemize}
          \item \impt{Núcleo Docente Estruturante}
          \item Pensa e atualiza o projeto pedagógico
          \item Indicado pelo colegiado
        \end{itemize}
      \end{block}
    \end{column}
  \end{columns}

  \vspace{0.3cm}
   
  \ntc{Reuniões: mínimo 2 por semestre, com relatórios de integralização curricular.}
\end{frame}

% ============================================================
\begin{frame}{O que o colegiado decide --- currículo}
  \begin{enumerate}
    \item Compatibilizar diretrizes dos componentes curriculares
    \item Regulamentar componentes para execução do PPC
    \item \impt{Deliberar sobre ementas e programas} propostos pelos departamentos
    \item Propor ao Conselho Superior o \impt{projeto pedagógico} e suas alterações
  \end{enumerate}

  \vspace{0.4cm}
   
  \ntc{Departamento propõe; colegiado aprova (ou pede ajuste) à luz do PPC.}
\end{frame}

% ============================================================
\begin{frame}{O que o colegiado decide --- trajetória do aluno}
  Questões \impt{individuais} do percurso acadêmico:
  \begin{itemize}
    \item Reopção de curso
    \item \impt{Equivalência de disciplinas}
    \item Desligamento e jubilamento
    \item \impt{Aproveitamento de estudos}
    \item Portador de Diploma (PDG)
    \item Transferência e reingresso
    \item Mobilidade acadêmica
  \end{itemize}

  \vspace{0.3cm}
   
  \begin{alertblock}{Heurística}
    ``Minha matéria $X$ feita na situação $Y$ vale para o meu curso?''\\
    $\Rightarrow$ \impt{Colegiado.} Não o departamento que oferece a disciplina.
  \end{alertblock}
\end{frame}

% ============================================================
\begin{frame}{O que o colegiado decide --- conclusão}
  \begin{itemize}
    \item \impt{Coordenar a orientação acadêmica} (integralização curricular)
    \item \impt{Indicar à PROGRAD os candidatos à colação de grau}\\
          \ntc{``Este aluno cumpriu tudo o que o PPC exige.''}
    \item Indicar membros do NDE
    \item Recomendar ao departamento: \impt{abertura de vagas e de turmas}
  \end{itemize}
\end{frame}

% ============================================================
\section{Departamento {\itshape vs.} Colegiado}

\begin{frame}{Comparação --- parte 1}
  \small
  \renewcommand{\arraystretch}{1.3}
  \begin{tabular}{@{}p{3.2cm} p{4.6cm} p{4.6cm}@{}}
    \toprule
    \textbf{Dimensão} & \textbf{Departamento} & \textbf{Colegiado} \\
    \midrule
    Organiza-se por & área de conhecimento & curso específico \\
    Lotação de & docentes e técnicos & ninguém é \emph{lotado} \\
    Atribui & encargos didáticos & componentes ao PPC \\
    Propõe & ementas e programas & projeto pedagógico \\
    \bottomrule
  \end{tabular}
\end{frame}

% ============================================================
\begin{frame}{Comparação --- parte 2}
  \small
  \renewcommand{\arraystretch}{1.3}
  \begin{tabular}{@{}p{3.2cm} p{4.6cm} p{4.6cm}@{}}
    \toprule
    \textbf{Dimensão} & \textbf{Departamento} & \textbf{Colegiado} \\
    \midrule
    Aprova & planos de trabalho docentes & equivalências, aproveitamentos \\
    Decide sobre & contratação e encargo & desligamento e jubilamento do aluno \\
    Presidência & chefe (2 anos) & coordenador (2 anos) \\
    Discentes & representação na assembleia & representação no colegiado \\
    Reuniões/sem. & $\geq$ 2 & $\geq$ 2 \\
    \bottomrule
  \end{tabular}
\end{frame}

% ============================================================
\section{Requerimentos do aluno}

\begin{frame}{Requerimentos ao departamento}
  Apenas \impt{dois tipos} --- ambos sobre oferta de disciplinas:
  \begin{enumerate}
    \item \impt{Abertura de turma}
    \item \impt{Abertura de vaga} em disciplina
  \end{enumerate}

  \vspace{0.4cm}
   
  \begin{block}{Lógica}
    Departamento é quem oferece a disciplina.\\
    \ntc{Disponibilizar disciplina = pedido ao departamento.}
  \end{block}
\end{frame}

% ============================================================
\begin{frame}{Requerimentos ao colegiado --- matrícula excepcional}
  \begin{enumerate}
    \item Quebra de pré-requisito
    \item Quebra de período consecutivo
    \item Matrícula com excesso de horas
    \item Excesso de horas em disciplina facultativa
  \end{enumerate}

  \vspace{0.4cm}
   
  \begin{alertblock}{Atenção}
    Nesses quatro casos, \impt{não é permitido trancamento parcial} da disciplina cuja matrícula foi feita por requerimento.
  \end{alertblock}
\end{frame}

% ============================================================
\begin{frame}{Requerimentos ao colegiado --- validação de estudos}
  \begin{itemize}
    \item \impt{Aproveitamento de estudos} \\
          \ntc{(Res. CEPE 7.325 --- histórico + programa analítico)}
    \item \impt{Extraordinário aproveitamento} \\
          \ntc{(Res. CEPE 1.986 --- CRE $\geq 7{,}0$, top 12\%)}
    \item \impt{Exame de nivelamento em língua estrangeira} \\
          \ntc{(Res. CEPE 4.672)}
  \end{itemize}
\end{frame}

% ============================================================
\begin{frame}{Requerimentos ao colegiado --- vínculo e trajetória}
  \begin{itemize}
    \item Reopção de curso / reopção de habilitação
    \item Reingresso
    \item Transferência (outra IES)
    \item Portador de Diploma de Graduação (PDG)
    \item Mobilidade acadêmica (ANDIFES, internacional)
  \end{itemize}

  \vspace{0.3cm}
   
  \ntc{Em vários casos há processo seletivo da PROGRAD para vagas residuais --- mas a decisão acadêmica final é do \impt{colegiado de destino}.}
\end{frame}

% ============================================================
\begin{frame}{Requerimentos ao colegiado --- suspensão e conclusão}
  \begin{columns}[T,onlytextwidth]
    \begin{column}{0.48\textwidth}
      \begin{block}{Suspensão}
        \begin{itemize}
          \item Afastamento especial (até 4 anos, \impt{uma vez})
          \item Trancamento fora de prazo
          \item Licença maternidade/paternidade
          \item RETEF
        \end{itemize}
      \end{block}
    \end{column}
    \begin{column}{0.48\textwidth}
      \begin{block}{Atividades e conclusão}
        \begin{itemize}
          \item AACC
          \item 2ª chamada (falecimento, 7 dias)
          \item Colação de grau
        \end{itemize}
      \end{block}
    \end{column}
  \end{columns}
\end{frame}

% ============================================================
\begin{frame}{Nem tudo vai a departamento ou colegiado}
  Outros destinos comuns:
  \begin{itemize}
    \item \impt{Seção de Ensino / minhaUFOP}: matrícula, cadastro, documentos
    \item \impt{PROGRAD}: cancelamento institucional, auxílio a eventos
    \item \impt{SISBIN}: carteira estudantil
    \item \impt{PRACE / PROPP / PROEx}: bolsas (alimentação, moradia, IC, extensão)
    \item \impt{Conselho da Unidade}: recursos contra decisões de departamento ou colegiado
  \end{itemize}
\end{frame}

% ============================================================
\begin{frame}{Heurística prática}
  \begin{block}{Para o aluno decidir a quem recorrer}
    \begin{itemize}
      \item \impt{Disponibilizar disciplina} (oferta) $\Rightarrow$ Departamento
      \item \impt{Qualquer outra coisa acadêmica} (regra de matrícula, validação, mudança de vínculo, suspensão, conclusão) $\Rightarrow$ Colegiado
      \item \impt{Administrativo} (documento, cadastro) $\Rightarrow$ Seção de Ensino
      \item \impt{Financeiro/assistencial} (bolsa, auxílio) $\Rightarrow$ Pró-reitoria competente
    \end{itemize}
  \end{block}
\end{frame}

% ============================================================
\section{Casos concretos}

\begin{frame}{Casos 1 e 2}
  \begin{alertblock}{Caso 1}
    Aluno cursou Cálculo I em outra universidade e quer aproveitá-la no BCC.
  \end{alertblock}
   
  $\Rightarrow$ \impt{Colegiado do BCC} (aproveitamento de estudos).

  \vspace{0.5cm}
   
  \begin{alertblock}{Caso 2}
    Alunos acham a ementa de PCC104 desatualizada e querem propor reforma.
  \end{alertblock}
   
  $\Rightarrow$ \impt{Colegiado primeiro}, que dialoga com o \impt{DECOM}.\\
  \ntc{Ementa: proposta pelo departamento, aprovada pelo colegiado.}
\end{frame}

% ============================================================
\begin{frame}{Casos 3 e 4}
  \begin{alertblock}{Caso 3}
    Aluno do BCC quer mudar para Engenharia de Computação (reopção).
  \end{alertblock}
   
  $\Rightarrow$ \impt{Colegiados do BCC e da Engenharia de Computação}.

  \vspace{0.5cm}
   
  \begin{alertblock}{Caso 4}
    Docente do DECOM quer afastar-se por um ano para pós-doutorado.
  \end{alertblock}
   
  $\Rightarrow$ \impt{Departamento (Assembleia do DECOM)}.\\
  \ntc{Plano de capacitação $\to$ Conselho da Unidade.}
\end{frame}

% ============================================================
\begin{frame}{Casos 5, 6 e 7}
  \begin{alertblock}{Caso 5}
    DECOM precisa decidir quem ministra PCC104 no próximo semestre.
  \end{alertblock}
   
  $\Rightarrow$ \impt{Departamento} (atribuição de encargos).

  \vspace{0.3cm}
   
  \begin{alertblock}{Caso 6}
    Aluno concluiu tudo e quer colar grau.
  \end{alertblock}
   
  $\Rightarrow$ \impt{Colegiado do curso} (indica à PROGRAD).

  \vspace{0.3cm}
   
  \begin{alertblock}{Caso 7}
    Aluno quer questionar decisão do coordenador de curso.
  \end{alertblock}
   
  $\Rightarrow$ \impt{Conselho da Unidade Acadêmica} (instância superior).
\end{frame}

% ============================================================
\section{Por que isso importa}

\begin{frame}{Duas portas de entrada}
  \begin{block}{Separação estrutural}
    Quem \impt{oferece a disciplina} $\neq$ quem \impt{coordena o curso}.
  \end{block}

  \vspace{0.4cm}
   
  \begin{itemize}
    \item Docente especializado é gerido junto com seus pares de área
    \item Projeto formativo de cada curso é pensado a partir da perspectiva específica do curso
    \item Várias contribuições departamentais convergem em um único PPC
  \end{itemize}
\end{frame}

% ============================================================
\begin{frame}{Representação discente: canais complementares}
  O corpo discente tem representação \impt{nas duas instâncias}:

  \vspace{0.3cm}
  \begin{columns}[T,onlytextwidth]
    \begin{column}{0.48\textwidth}
      \begin{block}{Assembleia de departamento}
        Decisões sobre:
        \begin{itemize}
          \item Oferta de disciplinas
          \item Corpo docente
        \end{itemize}
      \end{block}
    \end{column}
    \begin{column}{0.48\textwidth}
      \begin{block}{Colegiado de curso}
        Decisões sobre:
        \begin{itemize}
          \item Currículo
          \item Trajetória dos alunos
        \end{itemize}
      \end{block}
    \end{column}
  \end{columns}

  \vspace{0.3cm}
   
  \ntc{Atuação estratégica: saber em qual instância levantar cada questão.}
\end{frame}

% ============================================================
\begin{frame}{Como resolver problemas na UFOP}
  Ordem prática de consulta:
  \begin{enumerate}
    \item \impt{Colega de turma} --- mesma situação que você
    \item \impt{Veterano} --- já enfrentou algo parecido
    \item \impt{Página de orientações da PROGRAD} \\
          \ntc{\small\url{https://www.prograd.ufop.br/orientacoes-gerais-alunos-e-professores}}
    \item \impt{Colegiado do curso}
          \begin{itemize}
            \item CC: \ntc{\texttt{cocic@ufop.edu.br}}
            \item BIA: \ntc{\texttt{cobia@ufop.edu.br}}
          \end{itemize}
  \end{enumerate}
\end{frame}

% ============================================================
\begin{frame}{Mensagem final}
  \begin{block}{Para levar para casa}
    \begin{itemize}
      \item \impt{Departamento}: onde o professor é lotado --- responde por oferta e corpo docente.
      \item \impt{Colegiado}: onde o curso é coordenado --- responde por trajetória, currículo e validação de percurso.
      \item \impt{Conselho da Unidade}: instância de recurso.
    \end{itemize}
  \end{block}

  \vspace{0.5cm}
  \begin{center}
    \large \impt{Saber a quem recorrer é parte da formação universitária.}
  \end{center}
\end{frame}

% ============================================================
\begin{frame}{Fontes}
  \footnotesize
  \begin{itemize}
    \item Estatuto da UFOP --- Resolução CUNI nº 1.868, de 17/02/2017 (e alterações).
    \item Regimento Geral da UFOP --- Resolução CUNI nº 1.959, de 28/11/2017 (e alterações).
    \item PROGRAD --- Orientações Gerais a Alunos e Professores. \\
          \ntc{\url{https://www.prograd.ufop.br/orientacoes-gerais-alunos-e-professores}}
    \item Resoluções CEPE: 1.744, 1.986, 2.871, 4.672, 5.709, 7.325.
    \item Resolução CONGRAD 250/2025.
  \end{itemize}
\end{frame}

\end{document}